\documentclass[11pt,a4paper]{article}
\usepackage[margin=2.2cm]{geometry}
\usepackage[T1]{fontenc}
\usepackage{amsmath,amssymb,xcolor,enumitem,booktabs,graphicx,tabularx}
\usepackage[normalem]{ulem}
\setlength{\parskip}{0.6em}\setlength{\parindent}{0pt}
\definecolor{newc}{HTML}{1f6f43}
\newcommand{\new}[1]{\textcolor{newc}{#1}}
\newcommand{\old}[1]{\textcolor{gray}{\sout{#1}}}
\newcommand{\note}[1]{{\small\itshape [#1]}}
\newcommand{\lab}[1]{\textsf{\footnotesize\detokenize{#1}}}
\renewcommand{\ref}[1]{\lab{#1}}
\renewcommand{\eqref}[1]{(\lab{#1})}
\newcommand{\nref}[1]{(\lab{#1})}
\newcommand{\citep}[1]{[\lab{#1}]}
\newcommand{\citet}[1]{\lab{#1}}
\newcommand{\mockup}[1]{\fbox{\small mockup: #1}}
\title{Proposed change, ``Locks, larger fleets and a fleet reward test''}
\date{25 November 2026, version 4}
\begin{document}\sloppy
\maketitle

\section*{How to read this}
This follows the discussion of 25 November on the performance test at 8, 32
and 64 agents. Agreed so far:
\begin{itemize}[nosep]
  \item Why the decentralised controllers lock is a result and a discussion
    point, not part of the Method. The Method keeps the design facts with
    their reasons and loses the sentences that predict behaviour.
  \item The Discussion may not bring new information. So the run that tests
    the explanation is described in the Setup, reported in the Results and
    only read in the Discussion.
  \item The test gives both decentralised controllers a fleet
    reward like CRAMP's. It is a test of the explanation, not a new
    controller, and the main results do not change.
  \item New in version 4: the five questions are answered (see the end). A
    lock is now two agents on neighbouring vertices that both stand still,
    each holding its task, for 200 steps in a row. The threshold is high on
    purpose and the text says why. The weight gets its reason. The lock
    counts wait for the replays on the Lenovo, so they are shown as
    \mockup{waiting}. Both fleet reward runs are training and neither has
    collapsed.
  \item New in version 3: both decentralised controllers get the fleet
    reward, and the new runs are trained exactly as the long 16-agent runs,
    from the same starting policies for the same 2\,000\,000 steps. The long
    runs are then the comparison, and no separate control is trained. The new
    policies are played at 32 and 64 agents. Locks are described as part of
    every larger fleet, with shared credit making them shorter.
  \item New in version 2: why tasks and not orders, the noise of rewards
    from agents far away, the weight fixed before the run, the test added after
    the results but designed before it ran, one training run per arm, task
    against order in the Discussion, and what happens if run A collapses.
  \item COMA and Wolpert and Tumer are not cited. COMA's critic sees the whole
    warehouse, which the Method rejects. CRAMP and van Knippenberg et al.\
    are the precedent, and both are already cited.
\end{itemize}
Black text stays, \new{green} is new and \old{struck} text goes.
\note{Notes in brackets} say why. Nothing below is in the thesis yet. The
run has not been made, so its results are mockups. The questions are at the
end.

\section*{Part 1. What was checked}
\begin{itemize}[nosep]
  \item \textbf{Performance test} (\texttt{results/L\_results/scaling/perf/},
    240 of 240 episodes exit 0). The congestion-aware 16-agent policies at
    their [A111] checkpoints, 20 test seeds, 10\,000 steps, 5000 warm-up, at
    8, 32 and 64 agents. Main loads 0.09, 0.3, 0.6.
  \item \textbf{Token Passing at 32 and 64} was run in the Appendix~E load
    sweep (fixed storage, 20 seeds). Capacity 0.5 and 1.0, so the main load
    is 0.6 of capacity at every fleet size. At 32 agents it keeps up on 20 of
    20 seeds even at 0.6. The load is not too heavy. The learned controllers'
    full fleets (busy share 1.00) and backlogs of 1700 to 2100 are their own.
  \item \textbf{The 32-agent trained policies are old.} Trained on commit
    \texttt{0c5427b}, before [A93], [A96], [A111] and the observation change of
    \#212. Only decentralised fixed and demand and section-based fixed
    finished, none under congestion-aware storage and none with messages. They
    are left out and not mentioned. Every result at 8, 32 and 64 agents comes
    from the 16-agent policies. A backup with their settings, commits and
    training curves is in \texttt{thesis-progress/results/training/fleet\_32/}.
  \item \textbf{Credit} (\texttt{training/env.py}, \lab{eq:credit}). The
    credit group is the whole fleet for the centralised controller, the zone
    for the section-based one and the agent alone for the decentralised one.
    So only the decentralised controllers have no shared reward.
  \item \textbf{Reward in a lock} (\lab{eq:reward}, Appendix~B). A chosen
    wait costs $c_{\mathrm{wait}} = 0.2$. A refused move costs $0.2 +
    n_{\mathrm{override}} = 0.7$. Backing out costs a move and loses
    progress. van Knippenberg et al.\ set waiting above moving for the same
    reason, to keep agents moving.
  \item \textbf{CRAMP} (Pham and Bera, \texttt{papers/2309.10275v3.pdf},
    Section~III-D). Decentralised, GNN messages. ``To foster teamwork, agents
    receive a significant reward when all complete an episode together but
    face penalties for hindering others.'' $R_e = +20$ to all agents when all
    reach their goals. Their limitation: deadlocks when two robots enter one
    corridor from opposite ends, ``like prioritizing one robot's passage''.
  \item \textbf{Locks with the new definition} (two agents on neighbouring
    vertices both stand still, each holding its task, 200 steps in a row),
    seeds out of 20, congestion-aware storage. The old count (any two agents
    holding tasks 500 steps) grew with the fleet without a jam and did not
    match the text. The new counts come from \texttt{longest\_pair\_stand\_still}
    (PR \#216) in the replays on the Lenovo, 18 of 70 done. Only episodes
    where one agent stood still for 200 steps or more can lock, and those are
    the ones replayed. Token Passing's longest single stand-still is 23, 30,
    55 and 64 steps at 8, 16, 32 and 64 agents.
\end{itemize}

\begin{center}\small
\begin{tabular}{l cccc c cccc}
\toprule
& \multicolumn{4}{c}{served share (seeds that keep up)} && \multicolumn{4}{c}{seeds that lock} \\
\cmidrule(lr){2-5}\cmidrule(lr){7-10}
controller & 8 & 16 & 32 & 64 && 8 & 16 & 32 & 64 \\
\midrule
centralised   & 1.00 (18) & 0.81 (5) & 0.84 (2) & 0.94 (18) && ? & ? & ? & ? \\
section-based & 0.99 (18) & 0.79 (6) & 0.87 (7) & 0.94 (17) && ? & ? & ? & ? \\
decentralised & 0.97 (17) & 0.80 (6) & 0.79 (2) & 0.70 (3) && ? & ? & ? & ? \\
with messages & 0.99 (19) & 0.34 (0) & 0.51 (1) & 0.69 (10) && ? & ? & ? & ? \\
\bottomrule
\end{tabular}
\end{center}
\note{The lock columns wait for the replays. With the old count, nearly every
seed at 64 agents ``locked'' while the centralised and section-based
controllers still served 0.94, which is why the definition changed.}

\section*{Part 2. Method, \lab{sec:method:resolution} (4.Implementation)}
\note{Design facts stay, each with its reason. The sentences that say what an
agent will learn go to the Results and the Discussion.}

No rule is added to break the lock. A cap on the open tasks of a
workstation was rejected, because it would change the assignment rule that
every controller shares and works against orders whose lines all go to the
same workstation. Conflict resolution gets no rule for it either, such as
making the agent at the workstation back out, so that whether the learned
controllers learn to back out is part of what is measured. A decentralised
agent's credit is its own reward only (Section~\ref{sec:method:training})\old{,
so it gains nothing when it backs out for another agent}\new{.} The
section-based and centralised controllers share the credit of their
group\old{, so they have a reason to learn it}\new{.} The crowding
term is the same whether an agent waits or backs out. An agent whose move
is refused pays the wait cost and the override term of
\nref{eq:reward}, and an agent that chooses to wait pays the wait cost
alone. \old{So in front of a workstation, choosing to wait is cheaper than
asking for a vertex that stays taken.} In a circle every agent in it asks to
move and is refused, so all of them pay. \old{A message can tell a decentralised
agent that another agent is stuck, but it does not change the agent's
reward, so backing out still gains it nothing. Whether messages make the
lock rarer is part of RQ4.} Token Passing does not lock this way, for the
reason given in Section~\ref{sec:method:reference}.

\section*{Part 3. Setup, a new paragraph after ``Training'' (5.Experiments)}
\note{The test is part of the experiment, not of the method, so it is
described where the other runs are. The full design and the training curves
go to a new section of Appendix~L, \lab{app:results:fleetreward}.}

\new{\textbf{The fleet reward test.} The two decentralised controllers hold
tasks longest, and they are the only ones whose credit is their own reward
alone. One reading is that their reward gives an agent no reason to back
out. The test checks that reading under congestion-aware storage. Each
decentralised controller, with and without messages, is trained once more
exactly as its long 16-agent run, from the same starting policy and for the
same 2\,000\,000 steps, with a fleet reward added. The long runs are the
comparison, so the reward is the only difference. If the new runs lock less,
the shared reward was the missing piece. If they lock as often, it was not,
and the cause lies in training or elsewhere. The test was added after the main results showed the locks. Its design, its
measures and what each outcome means were written down before it ran
(Appendix~\ref{app:results:fleetreward}).}

\new{The fleet reward follows CRAMP \citep{pham2024optimizingcrowdawaremultiagentpath},
which gives every agent a bonus when the whole fleet reaches its goals.
Tasks here keep arriving, so there is no moment when all are done. Instead,
each time any agent finishes a task, every agent gets a share of its
delivery bonus. An agent then gains when another agent finishes, also when
it backed out to let it pass. The share is the mean over the fleet, so on
average an agent earns as much from the fleet's deliveries as from its own.
The test asks whether shared credit changes what the agents do, not which
weight works best, so the weight was set to this one value before the run
and not tuned.}

\new{The bonus follows finished tasks, not finished orders, for three reasons.
The agent's own reward already pays per task, so sharing the same bonus
changes only who is credited and nothing else. No controller sees which
order a task belongs to, so an order bonus would pay for something the agent
cannot see. A task also finishes every few steps, while an order waits for
its last line, and a frequent reward is easier to learn from. The same
problem remains in a smaller form. An agent is also paid when an agent far
away finishes a task, which it neither sees nor helped. Taking the mean over
the fleet keeps that share small, but it adds noise to what the agent learns
from, and the test shows whether the noise outweighs the gain.}

\new{Formally, with $m$ agents and weight $n_{\mathrm{fleet}}$, the reward of
agent $a_i$ in the test run is}
\new{\begin{equation}
  R^{\mathrm{fleet}}_i(t) = R_i(t) + n_{\mathrm{fleet}}\,\frac{1}{m}
  \sum_{j=1}^{m} n_{\mathrm{deliver}}\,\mathbb{1}[a_j \text{ completes a task at } t],
  \qquad n_{\mathrm{fleet}} = 1 .
  \label{eq:fleetreward}
\end{equation}}
\new{The agent still acts on its own window and its messages. Only its
reward reports on the whole fleet, so the test steps past the rule that a
controller learns only from what it decides for
(Section~\ref{sec:method:training}). It does so on purpose, to test the
explanation, and it is not a fourth controller. Both runs are evaluated as
the other 16-agent runs and then played at 32 and 64 agents. Each
controller is one training run. The 20 test seeds show how much a result varies between
warehouses, not between training runs, so a small difference between the two
runs is not taken as an effect.}

\note{$n_{\mathrm{fleet}}$ is new, add it to \lab{tab:notation}. The sum
includes $a_i$ itself, as CRAMP's team reward does.}

\section*{Part 4. Results (5.Experiments)}

\subsection*{4.0. ``RQ1'', where a lock is defined (before \lab{tab:exp:locks})}
\note{The text and the code now agree. The 16-agent counts in
\lab{tab:exp:locks} change with the replays.}

Locks explain most of the decentralised controller's failure at 16 agents
(Table~\ref{tab:exp:locks}). \old{A seed locks when two agents hold the same
task for 500 steps or more.} \new{A seed locks when two agents on
neighbouring vertices both stand still, each holding its task, for 200 steps
in a row. The threshold is set high on purpose. The measure only watches and
never stops a run, so two agents that wait and then clear the way themselves
do not count as a lock. A lock is a pair that stays stuck long after the
controller should have resolved it.} The decentralised controller locks on
\mockup{waiting} seeds\ldots

\subsection*{4a. A new paragraph after ``RQ3, load'', larger fleets}
\note{The step time paragraphs already say the 32 and 64 policies are the
16-agent ones. This reports how well they serve.}

\new{\textbf{Larger fleets.} The policies trained at 16 agents serve well at
8 agents, where every learned controller keeps up on 17 to 19 of 20 seeds
(Table~\ref{tab:exp:fleets}). At 32 agents all of them fall behind, and at 64
agents the centralised and section-based controllers recover to a served
share of 0.94 while the two decentralised ones stay near 0.7. Token Passing
keeps up at 32 and 64 agents at the same share of its capacity
(Appendix~\ref{app:load}), so the larger warehouses are not too heavy.}

\new{\mockup{Locks at each fleet size, from the replays. One sentence on how
the counts change with the fleet and one on how they follow the served
share. The version 3 sentences used the old count and are dropped.}}

\note{Table \lab{tab:exp:fleets}: the table of Part 1, served share with the
seeds that keep up, and the seeds that lock. Caption: ``Served share at each
fleet size of the policies trained at 16 agents, and the test seeds that
lock, out of 20.''}

\subsection*{4b. ``RQ4, communication'', the end}
The decentralised controller with messages still locks on 15 to 20 seeds,
so messages do not let two agents back out of a lock. \new{The same holds at
32 and 64 agents (Table~\ref{tab:exp:fleets}).}

\subsection*{4c. A new paragraph after RQ4, the fleet reward test}
\new{\textbf{The fleet reward test.} \mockup{seeds that lock and served share
and holds of 1000 steps at 16, 32 and 64 agents, each decentralised
controller with the fleet reward against its long run, paired by seed. One sentence per outcome, no reading.}}

\section*{Part 5. Discussion (5.Experiments)}

\subsection*{5a. A new paragraph after ``Why the learned controllers fall
behind''}
\note{This reads the results of Part 4 and adds nothing new. It points back
to the tragedy of the commons in the Background.}

\new{\textbf{Why the decentralised controllers lock.} Every learned
controller locks more as the fleet grows, because agents meet more often.
How long a lock lasts follows how the reward is shared. The more of the fleet
an agent's credit covers, the sooner its locks end. The centralised controller
shares the whole fleet's reward and its locks end soonest, the section-based
controller shares its zone's, and the decentralised controller has only its
own. A decentralised agent in a lock
pays to back out, by a move and lost progress, while the agent it lets pass
gains. Waiting is also cheaper than asking for a vertex that stays taken. So
for each agent alone, waiting is the better choice, and when both wait,
neither moves. This is the pattern of the tragedy of the commons from the
Background, choices that are sensible for each agent and bad for the fleet.}

\new{Messages were the local signal meant to solve this, and they did not.
A message can carry that another agent is stuck, but it does not change what
an agent is paid for. The messages are learned, so what they carry is not
known, and the runs cannot tell why the controller with messages locks more
often than the one without. Part of it may be the fleet size, since at 32
agents an agent receives more messages than in training.}

\new{The fleet reward pays for tasks because the controllers do not see
orders. A task bonus fits that, but it also pays an agent for tasks far away
that it neither saw nor helped, so part of what it learns from is noise. An
order bonus is closer to what the warehouse wants, since the customer waits
for the whole order. It only makes sense once the agent can see its order.
With deadlines on orders, as proposed below, the agent's features would
carry which order its task belongs to and the time left, and the fleet
reward could then follow finished orders, or orders finished in time.}

\new{\mockup{The fleet reward test, read. If the runs with the fleet reward
lock less than their long runs, the missing piece was the shared reward. If
they lock as often, it was not, and training or the lock itself is the
better explanation.}}

\subsection*{5b. ``What could have been done differently'', the item on locks}
\old{\emph{A way out of locks.} Most of the decentralised controller's loss is
locks. A shared credit, or a rule that makes one of two locked agents
give way, would let them back out.}

\new{\emph{A way out of locks.} \mockup{Depends on the test.} Three ways
follow from the reading above. A penalty for holding a vertex that another
agent's move is refused onto, as CRAMP penalises hindering others, keeps the
signal local. A credit group of the agents within an agent's own window, as
PRIMAL tried \citep{sartoretti2019primal}, shares the reward without reaching
past what the agent sees. A rule in conflict resolution that gives one of the
two agents priority, as CRAMP proposes for its own head-on locks, ends the
lock without learning. Training at the fleet size the controller will run at
would let it meet the locks of that fleet.}

\note{COMA is left out. Its critic sees the whole warehouse, which
\lab{sec:method:training} rejects with a reason, and offering it here would
contradict that choice. If wanted, one clause in ``Threats to validity''
can say that the decentralised controllers learn from less than the usual
centralised critic gives, and may be held back by it.}

\section*{Part 6. The runs}
\begin{itemize}[nosep]
  \item \textbf{Two runs, one per decentralised controller,} under
    congestion-aware storage. Each starts from the same policy as its long
    run (the Kaggle dataset \texttt{slap-init-f16-decentralised}):
    without messages \texttt{decentralised/congestion}
    \texttt{checkpoint\_000027}, with messages
    \texttt{decentralised\_comm/congestion} \texttt{checkpoint\_000013}.
  \item \textbf{Same settings as the long runs.} 16 agents, arrival rate 0.18,
    16\,000-step episodes, 3 runners, 2\,000\,000 steps (about 42
    iterations), checkpoints every 3 iterations. Only \eqref{eq:fleetreward}
    is added. The comparison is the long run of the same controller
    (\texttt{decentralised/congestion\_long} and
    \texttt{decentralised\_comm/congestion\_long\_part2}), already evaluated.
    The text says all runs were trained the same way (author).
  \item \textbf{Code.} One option in the training config, off by default. A
    small PR from main. Before launch, one iteration with the option off is
    checked to match commit \texttt{6dbe094}, which trained the long runs.
  \item \textbf{Where (as running).} Without messages on Kaggle's GPU to
    about iteration 7, then continued on the PC once the centralised demand
    evaluation ends. With messages on Colab (L4) to about iteration 24, then
    continued on the PC after the run without messages. At iteration 4 both
    are at or above their long runs (tasks per step 0.215 against 0.210
    without messages, 0.127 against 0.086 with messages), so neither has
    collapsed.
  \item \textbf{Evaluation.} As every 16-agent run: 5 validation seeds choose
    the checkpoint at $f_{\mathrm{up}} = 0.9$, then 20 test seeds, 10\,000
    steps, 5000 warm-up. Then the chosen checkpoint at 32 and 64 agents on the
    20 performance test seeds.
  \item \textbf{Measures.} Seeds that lock (200 steps, new definition),
    longest pair stand-still, served share, seeds that keep up, the override count, paired by seed
    against the long run.
  \item \textbf{If a run collapses.} The training log is checked now and then.
    If tasks per step fall well below the long run's at the same iteration
    and keep falling for four iterations, the run is stopped and the test is
    dropped from the chapters (author). See Question 5.
\end{itemize}

\section*{Questions, answered 25 November}
\begin{enumerate}
  \item \textbf{Lock threshold.} 200 steps under the new definition, kept
    high on purpose because the measure only watches and should count pairs
    that stay stuck, not pairs the controller resolves itself (author). The
    reason is in 4.0.
  \item \textbf{Weight.} $n_{\mathrm{fleet}} = 1$, agreed. The text gives the
    reason, the test is about shared credit, not about tuning a weight
    (Part 3).
  \item \textbf{Placement.} Agreed. Setup paragraph, Results table and
    paragraph, Discussion reading, details in a new Appendix~L section.
  \item \textbf{The old 32-agent runs.} Left out and not mentioned. A backup
    is kept in thesis-progress (Part 1).
  \item \textbf{A collapse.} Neither run has collapsed so far, so the test
    goes in. If one collapses later, the one sentence in Appendix~L stays as
    proposed.
\end{enumerate}

\section*{Still to do before applying}
\begin{itemize}[nosep]
  \item Fill the lock columns and 4.0, 4a from the replays (18 of 70 done).
  \item Fill 4c and the Discussion mockup from the fleet reward runs once
    they are evaluated.
  \item Apply to the thesis as [A143] with a GAPS.md entry, and add
    $n_{\mathrm{fleet}}$ to \lab{tab:notation}.
\end{itemize}

\end{document}
